\documentclass[10pt,a4paper]{amsart}
\usepackage[margin=19mm]{geometry}
\usepackage{amsmath,amssymb,mathtools}
\usepackage[scr=rsfs,cal=euler]{mathalfa}
\usepackage{xfrac}
\usepackage[unicode,hidelinks,bookmarks=false]{hyperref}
\newcommand{\ie}{i.e.\ }
\newcommand{\cf}{cf.\ }
\newcommand{\resp}{respectively\ }
\newcommand{\Subsection}{\subsection}
\providecommand{\nobreakdash}{\nobreak\hbox{-}}
\setlength{\emergencystretch}{2em}
\multlinegap0pt
\usepackage{tikz}
\usetikzlibrary{cd}
\usepackage{mathtools}
\usepackage{amssymb}
\usepackage{xcolor}
\newtheorem{theorem}{Theorem}
\newtheorem{lemma}{Lemma}
\newtheorem{proposition}{Proposition}
\newtheorem{corollary}{Corollary}
\newcommand{\Hom}{\operatorname{Hom}}
\newcommand{\Q}{\mathbb Q}
\newcommand{\Z}{\mathbb Z}
\DeclareMathOperator{\ord}{ord}
\title{Strong Weil Degree Divisibility at Higher Levels}
\author{Daeyeol Jeon, Yongjae Kwon}
\date{Source: arXiv:2608.06054v2 (CC BY 4.0)}
\begin{document}
\maketitle

\noindent\emph{Source and licence.} Strong Weil Degree Divisibility at Higher Levels. Version arXiv:2608.06054v2 (CC BY 4.0), distributed by arXiv under the Creative Commons Attribution 4.0 International licence, http://creativecommons.org/licenses/by/4.0/, as recorded in the arXiv licence metadata for this exact version and shown on the abstract page https://arxiv.org/abs/2608.06054v2. Full licence notice: \texttt{licenses/arxiv-2608.06054-CC-BY-4.0.txt}.\par\smallskip

\noindent\textit{Editorial scope.} Counted unit: the article's theorem thm:denominator (source0.tex lines 1096--1110). The article's complete original proof of it is delivered with it (source0.tex lines 1112--1187, one \texttt{proof} environment of 76 lines). No mathematics is written, repaired or re-derived: the statement and the proof are the article's own lines, in the article's own notation, and no retained line is substituted or re-flowed.  Retained uncounted as the source-local context the proof needs: lines 508--514; lines 584--592; lines 726--734; lines 811--820; lines 864--872; lines 899--912; lines 973--979; lines 1063--1071.  Nothing was omitted from any retained span, so the package carries no omission record.  No retained line contains a citation, so the wrapper carries no bibliography and no bibliography entry.  The retained lines contain the article's own diagram code, which the wrapper typesets with the standard tikz packages; no image file is delivered and the package declares no asset dependency.  Editorial formatting only, in addition: an AMS \texttt{amsart} preamble that restates the article's own declarations and macros used by the retained lines, namely the environments \texttt{theorem}, \texttt{lemma}, \texttt{proposition}, \texttt{corollary} on separate counters, together with the standard packages those lines need.  The retained environments print in this document as Theorem 1, Lemma 1, Proposition 1, Corollary 1 in order of appearance, whereas the article prints the same items with the numbering of its own full document.  The complete native source of the article as distributed in the arXiv e-print for this exact version is bundled unchanged as \texttt{sources/206957/source0.tex}.
\begin{theorem}\label{thm:rational-old-basis}
Let \(E/\Q\) be the strong Weil curve of conductor \(M\), and let \(E'/\Q\)
belong to its \(\Q\)-isogeny class. Let \(N\) be a multiple of \(M\), and let
\(u:E\to E'\) be a nonzero isogeny over \(\Q\). Then
the old homomorphisms \(\{\Phi_{r,u}:r\mid N/M\}\) form a \(\Q\)-basis of
\(\Hom_\Q(J_0(N),E')\otimes_\Z\Q\).
\end{theorem}

\begin{lemma}\label{lem:optimal-primitivity}
Let \(E/\Q\) be the strong Weil curve of conductor \(M\), and let \(E'\) be an
elliptic curve over \(\Q\) in the \(\Q\)-isogeny class of \(E\). Then composition
with \(\Phi_E\) induces the isomorphism
\(\Hom_\Q(E,E')\xrightarrow{\sim}\Hom_\Q(J_0(M),E')\),
\(v\mapsto v\circ\Phi_E\), of \(\Z\)-modules.
If \(u\) generates \(\Hom_\Q(E,E')\), then \(u\circ\Phi_E\) is primitive in
\(\Hom_\Q(J_0(M),E')\).
\end{lemma}

\begin{equation}\label{eq:D-factor}
        D_{u\circ\pi_E}(R)=
        \begin{cases}
        1, & R=1,\\[3pt]
        \displaystyle
        \prod_{\ell}\ell^{\ord_\ell(c_{u\circ\pi_E})\max\{1,\ord_\ell(R)\}},
        & R>1.
        \end{cases}
\end{equation}

\begin{lemma}\label{lem:two-cusp}
Let \((x_r)_{r\mid R}\) be a family of rational numbers, and suppose that
\(\omega=\sum_{r\mid R}x_r\eta_r\) is the pullback of a N\'eron differential
under a morphism from \(X_0(N)\) to an abelian variety over \(\Q\).
Then
\begin{equation}\label{eq:gcd-condition}
        \gcd(r,R/r)x_r\in\Z
        \qquad(r\mid R).
\end{equation}
\end{lemma}

\begin{proposition}\label{prop:remove}
Let \(\ell\) be a prime, write \(a=\ord_\ell(c_{u\circ\pi_E})\), and let \(\Z_{(\ell)}\) denote
the localization of \(\Z\) at \(\ell\). Suppose that an element
\(G\in\Hom_\Q(J_0(N),E')\otimes_\Z\Z_{(\ell)}\) has old-basis expansion
\[
        G=\sum_{r\mid R}x_r\Phi_{r,u}.
\]
Then \(\ell^a x_r\in\Z_{(\ell)}\) for every \(r\mid R\) with \(\ell\nmid r\).
\end{proposition}

\begin{proposition}\label{prop:degeneracy-degree-ramification}
Let \(N\) be a positive integer, and let \(\ell\) be a prime with
\(\ell^2\mid N\). Put \(N'=N/\ell\), and let
\(\delta_\ell=\iota_{\ell,N,N'}:X_0(N)\to X_0(N')\) be the degeneracy map
analytically induced by \(\tau\mapsto \ell\tau\).
Then \(\deg(\delta_\ell)=\ell\), and \(\delta_\ell\) is totally ramified at
\(\infty\). More
precisely, if \(q_N\) and \(q_{N'}\) denote the standard parameters at
\(\infty\), then
\begin{equation}\label{eq:q-total-ramification}
        \delta_\ell^*(q_{N'})=q_N^\ell,
\end{equation}
and \(e_\infty(\delta_\ell)=\ell=\deg(\delta_\ell)\).
\end{proposition}

\begin{proposition}\label{prop:connected-degeneracy-kernel}
With the notation of Proposition
\ref{prop:degeneracy-degree-ramification},
\((\delta_\ell)_*:J_0(N)\to J_0(N')\) is surjective and
\(K_\ell:=\ker((\delta_\ell)_*)\) is geometrically
connected.
\end{proposition}

\begin{corollary}\label{cor:connected-local}
With the notation of Proposition \ref{prop:connected-degeneracy-kernel}, let
\(A/\Q\) be an abelian variety. Precomposition with \((\delta_\ell)_*\) has
saturated image after localization at any prime \(p\). More precisely, let
\(H\in\Hom_\Q(J_0(N'),A)\otimes_\Z\Q\). If
\(H\circ(\delta_\ell)_*\in
\Hom_\Q(J_0(N),A)\otimes_\Z\Z_{(p)}\),
then \(H\in\Hom_\Q(J_0(N'),A)\otimes_\Z\Z_{(p)}\).
\end{corollary}

\begin{theorem}\label{thm:denominator}
Let \(D_{u\circ\pi_E}(R)\) be defined by \eqref{eq:D-factor}. If
\[
        G=\sum_{r\mid R}x_r\Phi_{r,u}\in\Hom_\Q(J_0(N),E')
\]
is the expansion of \(G\) in the \(\Q\)-basis of Theorem
\ref{thm:rational-old-basis}, then
\begin{equation}\label{eq:global-denominator}
        D_{u\circ\pi_E}(R)x_r\in\Z
        \qquad(r\mid R).
\end{equation}
Equivalently,
\(D_{u\circ\pi_E}(R)\Hom_\Q(J_0(N),E')\subseteq
\bigoplus_{r\mid R}\Z\Phi_{r,u}\).
\end{theorem}

\begin{proof}
The case \(R=1\) is Lemma \ref{lem:optimal-primitivity} applied with
\(E'\) and \(u\), because then
\(\Phi_{1,u}=u\circ\Phi_E\) generates \(\Hom_\Q(J_0(M),E')\). Hence
\(x_1\in\Z\).
Assume \(R>1\) and fix a prime \(\ell\). Put
\(a=\ord_\ell(c_{u\circ\pi_E})\) and \(m=\ord_\ell(R)\). We prove
\begin{equation}\label{eq:local-denominator}
        \ell^{a\max\{1,m\}}x_r\in\Z_{(\ell)}
        \qquad(r\mid R)
\end{equation}
by induction on \(m\).

If \(m\le1\), Lemma \ref{lem:two-cusp}, applied to the pullback
\[
        (G\circ j_N)^*\omega_{E'}
        =
        c_{u\circ\pi_E}\sum_{r\mid R}x_r\eta_r ,
\]
implies \(\gcd(r,R/r)c_{u\circ\pi_E}x_r\in\Z\) for every \(r\mid R\). When
\(m\le1\), both \(\gcd(r,R/r)\) and
\(c_{u\circ\pi_E}/\ell^a\) are units in \(\Z_{(\ell)}\). It follows that
\(\ell^ax_r\in\Z_{(\ell)}\), which is \eqref{eq:local-denominator} in this
case.

Assume that \(m\ge2\). Proposition \ref{prop:remove} implies
\(\ell^ax_r\in\Z_{(\ell)}\) whenever \(\ell\nmid r\).
Subtracting the terms with \(\ell\nmid r\) from \(\ell^aG\)
leaves only those indexed by multiples of \(\ell\), which factor through
\(J_0(N/\ell)\). Set
\[
        G_1=
        \ell^aG-
        \sum_{\substack{r\mid R\\ \ell\nmid r}}\ell^ax_r\Phi_{r,u,N,M}
        \in\Hom_\Q(J_0(N),E')\otimes_\Z\Z_{(\ell)}.
\]
Put \(N'=N/\ell\), \(R'=R/\ell\),
\(\delta_\ell=\iota_{\ell,N,N'}\).
For each \(e\mid R'\), the degeneracy maps compose to give
\[
        \Phi_{\ell e,u,N,M}
        =
        \Phi_{e,u,N',M}\circ(\delta_\ell)_*.
\]
Equivalently, for each \(e\mid R'\), the diagram
\[
\begin{tikzcd}[column sep=large,row sep=large]
J_0(N) \arrow[r,"(\delta_\ell)_*"] \arrow[dr,swap,"\Phi_{\ell e,u,N,M}"] &
J_0(N') \arrow[d,"\Phi_{e,u,N',M}"] \\
& E'
\end{tikzcd}
\]
commutes.
Collecting the terms with index \(\ell e\), we obtain
\[
        G_1=H\circ(\delta_\ell)_*,
        \qquad
        H=\sum_{e\mid R'}\ell^ax_{\ell e}\Phi_{e,u,N',M}
        \quad\text{in }\Hom_\Q(J_0(N'),E')\otimes_\Z\Q.
\]
Corollary \ref{cor:connected-local} shows that
\(H\in\Hom_\Q(J_0(N'),E')\otimes_\Z\Z_{(\ell)}\).
Choose an integer \(t\) prime to \(\ell\) such that
\(tH\in\Hom_\Q(J_0(N'),E')\). Apply the induction hypothesis at level \(N'\)
to \(tH\). Since \(t\) is a unit in \(\Z_{(\ell)}\) and
\(\ord_\ell(R')=m-1\), we have
\[
        \ell^{a(m-1)}\bigl(\ell^ax_{\ell e}\bigr)
        =\ell^{am}x_{\ell e}\in\Z_{(\ell)}.
\]
For \(\ell\nmid r\), Proposition \ref{prop:remove} already proves
\(\ell^ax_r\in\Z_{(\ell)}\), which is stronger than the required bound because
\(m\ge2\). This completes the induction and proves
\eqref{eq:local-denominator}. Combining these local statements over all
primes \(\ell\) proves \eqref{eq:global-denominator}.
\end{proof}

\end{document}
